\chapter{电子自旋和原子能级的精细结构}

\problem{3.1}{$\mu^-$ 粒子被氢原子俘获形成 $\mu$ 氢原子，基态时其轨道运动的磁矩有多大？基态电子偶素中电子和正电子绕共同质心运动的磁矩有多大？（已知 $m_\mu=207m_e$。）}
% layout
\answerarea

\problem{3.2}{一束基态氢原子束按斯特恩--盖拉赫实验方法通过极不均匀的磁场，磁场梯度 $\partial B_z/\partial z=1.5\times10^2\ \mathrm{T\,m^{-1}}$，氢原子的速度为 $v=10^4\ \mathrm{m\,s^{-1}}$，磁场区的长度 $d_1=20\ \mathrm{cm}$，从磁铁到屏的漂移距离 $d_2=10\ \mathrm{cm}$（如图所示），求氢原子束在屏上的裂距 $\Delta z$。（已知氢原子的质量为 $1.6\times10^{-27}\ \mathrm{kg}$。）
\fig{3.2}
\sourcenote{截图中磁场梯度的单位印作 $\mathrm{T^2\,m^{-1}}$，这里按 $\partial B_z/\partial z$ 的量纲整理为 $\mathrm{T\,m^{-1}}$，数值未改。}}
% layout
\answerarea
\newpage

\problem{3.3}{一束自旋为 $1/2$、磁矩为 $\mu$ 的中性粒子束沿 $x$ 方向通过斯特恩--盖拉赫实验装置，结果该中性粒子束按粒子的 $\mu_z$ 值不同而分裂。如果入射束中的粒子磁矩是下列情况之一，粒子束通过斯特恩--盖拉赫实验装置后的行为如何？
\begin{enumerate}[label=（\alph*）]
\item 磁矩沿 $+z$ 方向极化（排列）；
\item 沿 $-z$ 方向极化；
\item 沿 $+y$ 方向极化；
\item 不极化。
\end{enumerate}}
% layout
\answerarea

\problem{3.4}{设电子是一个半径为 $10^{-17}\ \mathrm m$ 的经典小球，试由其自旋角动量的大小估算其赤道上一点的线速度。（已知 $m_e=9.1\times10^{-31}\ \mathrm{kg}$，$\hbar=1.05\times10^{-34}\ \mathrm{J\,s}$。）}
% layout
\answerarea
\newpage

\problem{3.5}{氢原子的 $2\mathrm p$ 态电子由于自旋--轨道相互作用产生的能级裂距等于电子自旋磁矩在磁感应强度 $B$ 为多少的磁场中产生的能级裂距？}
% layout
\answerarea

\problem{3.6}{对于量子数 $L=1$、$S=1/2$ 的原子，
\begin{enumerate}
\item 计算 $\boldsymbol L\cdot\boldsymbol S$ 可能的值；
\item 计算其总角动量可能的值。
\end{enumerate}}
% layout
\answerarea
\newpage

\problem{3.7}{今测得氢原子 $n=3$ 是由谱项自下至上依次相差 $0.1082\ \mathrm{cm^{-1}}$ 与 $0.0361\ \mathrm{cm^{-1}}$ 的三能级所构成的。试问其赖曼系的第二条实际是由波长相差多少的两谱线构成的？}
% layout
\answerarea

\problem{3.8}{光谱仪的分辨本领 $R=\lambda/\delta\lambda$ 为多大时，可以分辨出氢原子 $\mathrm H_\alpha$ 线（巴耳末系的第一条）的全部精细结构成分？（不考虑兰姆位移。）}
% layout
\answerarea
\newpage

\problem{3.9}{考虑氢原子的基态和 $n=2$ 的各态。
\begin{enumerate}
\item 标出图中四个态的原子态符号。
\item 形成氢原子能级分裂的物理原因有：兰姆移位、自旋--轨道耦合、相对论效应、超精细结构。
\begin{enumerate}[label=（\alph*）]
\item 对 $n=1$ 的态，要考虑哪些修正因素？
\item 对 $n=2,l=0$ 的态和 $n=2,l=1$ 的态要考虑哪些修正因素？
\item 按大小列出四个修正项的次序或指出哪些修正项数量级相同。
\end{enumerate}
\end{enumerate}
\fig{3.9}}
% layout
\answerarea
\newpage

\problem{3.10}{已知铯原子光谱的锐线系的第一条谱线双线结构的波数为 $\bar\nu_1=6805\ \mathrm{cm^{-1}}$、$\bar\nu_2=7359\ \mathrm{cm^{-1}}$。
\begin{enumerate}
\item 这两条谱线是由哪两个能级劈裂引起的？
\item 此两子能级的能量差为多少电子伏？
\item 若认为能级差是由电子自旋磁矩与电子轨道运动在电子处产生的磁场间的作用引起的，试估算该内部磁场 $B$ 的大小。（$\mu_B=5.79\times10^{-5}\ \mathrm{eV\,T^{-1}}$。）
\end{enumerate}}
% layout
\answerarea

\newpage

\problem{3.11}{已知钾原子的基态为 $4^2\mathrm S_{1/2}$，其电离能为 $4.32\ \mathrm{eV}$，主线系第一条谱线由波长为 $\lambda_1=769.90\ \mathrm{nm}$ 和 $\lambda_2=766.41\ \mathrm{nm}$ 的双线组成；漫线系的第一条谱线的精细结构为 $\lambda_3=1168.98\ \mathrm{nm}$、$\lambda_4=1177.17\ \mathrm{nm}$ 和 $\lambda_5=1177.14\ \mathrm{nm}$。
\begin{enumerate}
\item 试画出产生这些谱线的能级跃迁图；
\item 求 $4\mathrm S$、$4\mathrm P$、$3\mathrm D$ 各能态的能量及相应的有效电荷数 $Z^*$ 和量子数亏损。
\end{enumerate}}
% layout
\answerarea
\newpage

\problem{3.12}{已知钠原子的共振线波长为 $589.3\ \mathrm{nm}$，漫线系第一条谱线的波长为 $819.3\ \mathrm{nm}$，基线系第一条谱线波长为 $184.59\ \mathrm{nm}$，主线系的系限波长为 $241.3\ \mathrm{nm}$。试求相关的 $\mathrm S$、$\mathrm P$、$\mathrm D$ 和 $\mathrm F$ 能级的量子数亏损值。}
% layout
\answerarea

\problem{3.13}{原子在热平衡条件下处在不同能量状态的原子数目按玻尔兹曼分布：
\[
N=N_0\frac{g}{g_0}\exp\!\left[-\frac{E-E_0}{kT}\right].
\]
其中 $N_0$ 和 $N$ 分别是处在 $E_0$ 和 $E$ 状态的原子数，$g_0$ 和 $g$ 是相应能量状态的统计权重，$k=8.62\times10^{-5}\ \mathrm{eV\,K^{-1}}$ 是玻尔兹曼常数。从高温铯原子气体光谱中测出主线系第一条谱线的双线 $\lambda_1=894.35\ \mathrm{nm}$、$\lambda_2=852.11\ \mathrm{nm}$ 的强度比 $I_1:I_2=2:3$，试估算气体的温度。（已知相应的能级的统计权重 $g_1=2$、$g_2=4$。）}
% layout
\answerarea
\newpage

\problem{3.14}{某碱金属原子处于弱磁场中。
\begin{enumerate}
\item $^2\mathrm P_{3/2}$ 和 $^2\mathrm D_{5/2}$ 态的 $g$ 因子各为多少？
\item $^2\mathrm P_{3/2}$ 和 $^2\mathrm D_{5/2}$ 的磁分裂能级的相邻能级间隔哪个大？
\item $^2\mathrm P_{3/2}$ 和 $^2\mathrm D_{5/2}$ 态相距最远的能级间隔分别为多少？（用 $\mu_BB$ 来表示。）
\end{enumerate}}
% layout
\answerarea

\problem{3.15}{跃迁 $^2\mathrm P_{1/2}\to{}^2\mathrm S_{1/2}$ 在弱磁场中将分裂为几条谱线？其与原谱线的波数差为多少（用 $\sigma$ 为单位表示）？画出其能级跃迁图。钠的 $589.59\ \mathrm{nm}$ 谱线刚好对应此跃迁。要使其 $\sigma$ 偏振线与 $\pi$ 偏振线的波长差为 $0.600\ \mathrm{nm}$，则外磁场的磁感应强度应为何值（取三位有效数字）？}
% layout
\answerarea
\newpage

\problem{3.16}{对于 Na 原子（$Z=11$），
\begin{enumerate}
\item 考虑了相对论效应和自旋--轨道耦合后，量子数 $n=3$ 的能级数是多少？
\item 在强磁场下，量子数 $n=3$ 的能级数为多少？定性画出其能级分裂图。
\end{enumerate}}
% layout
\answerarea

\problem{3.17}{试问基态氢原子在 $1.00\ \mathrm T$ 的磁场中塞曼劈裂的能量差是多少？若要使电子的自旋变换方向，需要外加振荡电磁场的频率是多少？}
% layout
\answerarea
\newpage

\problem{3.18}{氢的 $1420\ \mathrm{MHz}$ 发射线是在射电天文学中发现最早、最重要的谱线之一。利用这条谱线可探测宇宙空间中氢原子的分布情况，它来自于氢原子基态的超精细结构间的跃迁。试计算氢原子基态超精细作用系数 $a$。（已知氢原子核自旋量子数 $I=1/2$；$a>0$。）}
% layout
\answerarea

\problem{3.19}{在某种碱金属原子的电子顺磁共振试验中，固定微波频率为 $2.0\times10^{10}\ \mathrm{Hz}$，当磁场强度 $B$ 调到 $1.787\ \mathrm T$ 时发生强烈的共振吸收，试计算碱金属原子该状态的朗德 $g$ 因子，并指出原子处在何种状态？}
% layout
\answerarea
\newpage

